% 第 14 题双散点图：按固定源内容流矢量层精确值复刻（1 pt → 0.03515 cm，与全库同制）。
% 原点取左幅 O（源页 (815.1758, 206.5898) pt），y 向上；右幅 O 在 x 偏移 +176.2461 pt。
% 坐标换算直接写在坐标表达式里（保留源页坐标可审计）：x_m = 页 x − O_x，y_m = 206.5898 − 页 y。
% 源实测：两幅各「左轴+下轴+右上直角折线框」130.914 × 130.914 pt（刻度端值 267，1 名次单位 = 0.490315 pt）；
% 散点 72 枚（左幅 35、右幅 37——源图与题干「37 位学生」不完全对应，照源复刻），字库 • 墨迹 2.45 pt；
% 标注：甲（左幅 (39.7,81.2) 点右）、乙（左幅 (126.6,67.5) 点右）、丙（右幅 (169.3,172.8) 点右）。
\begin{tikzpicture}[x=0.03515cm,y=0.03515cm,>=stealth,every node/.style={font=\small,inner sep=0pt},
  gl/.style={line width=0.398bp,line cap=round,line join=round}]
% ===== 左幅（语文）=====
\draw[gl,->] (0,0)--(149.6133,0);
\draw[gl,->] (0,0)--(0,149.6171);
\draw[gl] (130.9140,0)--(130.9140,130.9140)--(0,130.9140);
\node at (-11.2700,130.2400) {267};
\node at (131.2700,-7.8100) {267};
\node at (-7.6400,-8.2200) {\(O\)};
\node at (69.5300,-9.0500) {总成绩年级名次};
\node at (-8.6458,107.6200) {语}; \node at (-8.6458,97.0100) {文};
\node at (-8.6458,86.2700) {成}; \node at (-8.6458,75.4000) {绩};
\node at (-8.6458,64.5300) {年}; \node at (-8.6458,53.7400) {级};
\node at (-8.6458,42.9300) {名}; \node at (-8.6458,32.2200) {次};
\node at (28.7132,40.2998) {甲};
\node at (71.3702,33.4908) {乙};
\foreach \x/\y in {
816.2932/199.6777, 817.4135/154.7919, 817.4147/201.9215, 818.9112/204.9131,
819.2836/150.3038, 820.7801/101.3035, 823.7729/183.2193, 824.8945/134.2191,
826.0173/188.0817, 826.7646/127.1129, 830.8791/90.4566, 833.1234/145.0666,
834.6198/166.7616, 835.7414/150.3038, 839.1079/138.3335, 843.5978/187.7080,
844.7182/143.9450, 847.3380/182.0966, 848.4596/179.4785, 849.5812/97.1891,
851.0788/198.5549, 854.8172/141.3258, 860.4281/122.9985, 873.8942/147.3109,
876.1385/155.1661, 877.2595/173.4934, 877.6344/109.1593, 878.0080/130.8531,
889.9777/192.5698, 891.0999/141.3264, 896.7102/97.9388, 906.4355/102.0532,
906.8092/178.7300, 909.0535/203.0430, 913.1679/140.2048}
  \fill ({\x-815.1758},{206.5898-\y}) circle[radius=1.225];
% ===== 右幅（数学）=====
\draw[gl,->] (176.2461,0)--(325.8594,0);
\draw[gl,->] (176.2461,0)--(176.2461,149.6171);
\draw[gl] (307.1601,0)--(307.1601,130.9140)--(176.2461,130.9140);
\node at (165.0761,130.2400) {267};
\node at (307.5161,-7.8100) {267};
\node at (169.7061,-8.2200) {\(O\)};
\node at (246.8761,-9.0500) {总成绩年级名次};
\node at (167.7003,107.6200) {数}; \node at (167.7003,97.0100) {学};
\node at (167.7003,86.2700) {成}; \node at (167.7003,75.4000) {绩};
\node at (167.7003,64.5300) {年}; \node at (167.7003,53.7400) {级};
\node at (167.7003,42.9300) {名}; \node at (167.7003,32.2200) {次};
\node at (268.6382,85.2698) {丙};
\foreach \x/\y in {
992.9134/183.2189, 993.2883/194.0658, 993.6625/201.9211, 994.7841/204.1654,
995.9063/189.9514, 996.6542/200.4246, 997.7764/179.8524, 1000.3944/164.5168,
1002.2647/179.8524, 1002.2651/203.0426, 1003.3872/189.9514, 1008.6233/183.2189,
1010.4934/176.1116, 1011.9898/178.3560, 1013.4856/140.2038, 1016.1042/176.4853,
1019.8438/183.2189, 1022.8367/161.5239, 1023.5846/164.5168, 1026.2026/116.6399,
1026.2026/146.1883, 1026.5775/153.6687, 1028.4476/148.8063, 1031.0656/183.2189,
1038.1723/129.3569, 1052.3857/112.5255, 1053.8822/201.9211, 1055.0037/106.9146,
1055.0037/117.7615, 1056.5002/97.5629, 1068.8435/81.4794, 1069.2172/100.9294,
1074.4544/121.8753, 1082.3084/196.3102, 1084.5534/94.5706, 1087.9193/105.4175,
1092.0337/87.8382}
  \fill ({\x-991.4219+176.2461},{206.5898-\y}) circle[radius=1.225];
\end{tikzpicture}
